\documentclass{beamer}
\usetheme{Madrid} % You can change this theme if you like

\title[Labor Market and Unemployment]{Labor Market and Unemployment}
\subtitle{Contemporary Economic Issues: ECON 204}
\author{Muhebullah Karimzada}
\date{Winter 2026}
\begin{document}

% Title Slide
\begin{frame}
    \titlepage
\end{frame}



% Section 1: What is the Labor Market?
\section{What is the Labor Market?}
\begin{frame}{What is the Labor Market?}
    \begin{itemize}
        \item The labor market is where workers supply labor and employers demand it.
        \item \textbf{Key Players:}
        \begin{itemize}
            \item \textbf{Workers:} Supply labor in exchange for wages.
            \item \textbf{Employers:} Demand labor to produce goods/services.
            \item \textbf{Government:} Influences through policies like minimum wages, taxes, etc.
        \end{itemize}
        \item \textbf{Key Concepts:}
        \begin{itemize}
            \item Wages: Price of labor.
            \item Employment: Quantity of labor exchanged.
            \item Unemployment: Workers who cannot find jobs.
            \item Human Capital: Skills, knowledge, and experience of workers.
        \end{itemize}
    \end{itemize}
\end{frame}

% Section 2: Understanding Unemployment
\section{Understanding Unemployment}
\begin{frame}{Understanding Unemployment}
    \textbf{Definition:} Unemployment occurs when people who are willing and able to work cannot find jobs.
    \vspace{0.5cm}
    \textbf{Types of Unemployment:}
    \begin{itemize}
        \item \textbf{Frictional:} Short-term, occurs during job transitions.
        \item \textbf{Structural:} Long-term, mismatch between skills and jobs.
        \item \textbf{Cyclical:} Due to economic downturns or recessions.
        \item \textbf{Seasonal:} Industries with seasonal demand (e.g., agriculture).
        \item \textbf{Hidden/Underemployment:} Working below skill level or fewer hours.
    \end{itemize}
\end{frame}

% Section 3: Measuring Unemployment
\section{Measuring Unemployment}
\begin{frame}{Measuring Unemployment}
    \textbf{Key Metrics:}
    \begin{itemize}
        \item \textbf{Unemployment Rate:}
        \[
        \text{Unemployment Rate} = \frac{\text{Number of Unemployed}}{\text{Labor Force}} \times 100
        \]
        \item \textbf{Labor Force Participation Rate:}
        \[
        \text{Participation Rate} = \frac{\text{Labor Force}}{\text{Working-Age Population}} \times 100
        \]
        \item \textbf{Discouraged Workers:} Individuals who stop looking for jobs are excluded from the labor force.
    \end{itemize}
\end{frame}

% Section 4: Causes of Unemployment
\section{Causes of Unemployment}
\begin{frame}{Causes of Unemployment}
    \begin{itemize}
        \item \textbf{Macroeconomic Shocks:} Recessions, financial crises, or pandemics reduce demand.
        \item \textbf{Technological Change:} Automation displaces workers.
        \item \textbf{Globalization:} Jobs outsourced to lower-cost countries.
        \item \textbf{Rigid Labor Markets:} Regulations may discourage hiring.
        \item \textbf{Mismatches:} Skills or geographic mismatches between workers and jobs.
    \end{itemize}
\end{frame}

% Section 5: Consequences of Unemployment
\section{Consequences of Unemployment}
\begin{frame}{Consequences of Unemployment}
    \textbf{Economic Costs:}
    \begin{itemize}
        \item Loss of income for individuals.
        \item Lower tax revenues and higher government spending on welfare.
        \item Lost potential output.
    \end{itemize}
    \textbf{Social Costs:}
    \begin{itemize}
        \item Increased poverty and inequality.
        \item Mental health challenges (stress, anxiety, depression).
        \item Social unrest and political instability.
    \end{itemize}
\end{frame}

% Section 6: Policies to Address Unemployment
\section{Policies to Address Unemployment}
\begin{frame}{Policies to Address Unemployment}
    \textbf{Demand-Side Policies:}
    \begin{itemize}
        \item \textbf{Fiscal Policy:} Government spending and tax cuts to boost demand.
        \item \textbf{Monetary Policy:} Lowering interest rates to encourage investment.
    \end{itemize}
    \textbf{Supply-Side Policies:}
    \begin{itemize}
        \item Education and training programs.
        \item Labor market flexibility (e.g., reducing barriers to hiring/firing).
        \item Tax incentives for job creation.
    \end{itemize}
    \textbf{Safety Nets:}
    \begin{itemize}
        \item Unemployment insurance.
        \item Public works programs for temporary jobs.
    \end{itemize}
\end{frame}

% Section 7: Discussion Questions
\section{Discussion Questions}
\begin{frame}{Discussion Questions}
    \begin{itemize}
        \item How does automation impact different types of unemployment?
        \item What policies should governments prioritize during a recession?
      
    \end{itemize}
\end{frame}


% Section 1: Numerical Problems
\section{Numerical Problems}

% Problem 1: Calculating the Unemployment Rate
\begin{frame}{Problem 1: Calculating the Unemployment Rate}
The labor force in a country consists of 150 million people, of which 10 million are unemployed.
\begin{enumerate}
    \item Calculate the unemployment rate.  
    \item If 2 million unemployed people stop looking for jobs, what happens to the unemployment rate?  
\end{enumerate}
\end{frame}



% Section 1: Numerical Problems
\section{Numerical Problems}

% Problem 1: Calculating the Unemployment Rate
\subsection{Problem 1: Calculating the Unemployment Rate}
\begin{frame}{Problem 1: Calculating the Unemployment Rate}
The labor force in a country consists of 150 million people, of which 10 million are unemployed.
\begin{enumerate}
    \item Calculate the unemployment rate.  
    \item If 2 million unemployed people stop looking for jobs, what happens to the unemployment rate?  
\end{enumerate}
\end{frame}

\begin{frame}{Answer to Problem 1}
(a) Unemployment Rate:
\[
\text{Unemployment Rate} = \frac{\text{Unemployed}}{\text{Labor Force}} \times 100 = \frac{10}{150} \times 100 = 6.67\%
\]

(b) Adjusted Unemployment Rate:  
If 2 million unemployed stop looking for jobs, the labor force becomes \(150 - 2 = 148 \, \text{million}\), and the unemployment rate becomes:
\[
\text{Unemployment Rate} = \frac{8}{148} \times 100 \approx 5.41\%
\]
\end{frame}

% Problem 2: Labor Force Participation Rate
\subsection{Problem 2: Labor Force Participation Rate}
\begin{frame}{Problem 2: Labor Force Participation Rate}
The working-age population of a country is 200 million. Of these, 140 million are in the labor force, and 60 million are not participating in the labor market.
\begin{enumerate}
    \item Calculate the labor force participation rate.  
\end{enumerate}
\end{frame}

\begin{frame}{Answer to Problem 2}
\[
\text{Labor Force Participation Rate} = \frac{\text{Labor Force}}{\text{Working-Age Population}} \times 100
\]
Substituting values:
\[
\text{Labor Force Participation Rate} = \frac{140}{200} \times 100 = 70\%
\]
\end{frame}

% Problem 3: Employment Rate
\subsection{Problem 3: Employment Rate}
\begin{frame}{Problem 3: Employment Rate}
In a small country, the labor force is 25 million, and 22 million people are employed.
\begin{enumerate}
    \item Calculate the unemployment rate.  
    \item Calculate the employment rate.
\end{enumerate}
\end{frame}

\begin{frame}{Answer to Problem 3}
(a) Unemployment Rate:
\[
\text{Unemployed} = 25 - 22 = 3 \, \text{million}
\]
\[
\text{Unemployment Rate} = \frac{3}{25} \times 100 = 12\%
\]

(b) Employment Rate:
\[
\text{Employment Rate} = \frac{\text{Employed}}{\text{Labor Force}} \times 100 = \frac{22}{25} \times 100 = 88\%
\]
\end{frame}

% Problem 4: Cyclical Unemployment
\subsection{Problem 4: Cyclical Unemployment}
\begin{frame}{Problem 4: Cyclical Unemployment}
The natural rate of unemployment in a country is estimated to be 4\%. During a recession, the actual unemployment rate rises to 8\%.
\begin{enumerate}
    \item What is the cyclical unemployment rate?  
    \item If the labor force is 50 million, how many people are cyclically unemployed?  
\end{enumerate}
\end{frame}

\begin{frame}{Answer to Problem 4}
(a) Cyclical Unemployment Rate:
\[
\text{Cyclical Unemployment Rate} = \text{Actual Unemployment Rate} - \text{Natural Unemployment Rate}
\]
\[
\text{Cyclical Unemployment Rate} = 8\% - 4\% = 4\%
\]

(b) Number of Cyclically Unemployed:
\[
\text{Cyclically Unemployed} = 4\% \times 50 \, \text{million} = 0.04 \times 50 = 2 \, \text{million}
\]
\end{frame}

% Problem 5: Structural vs. Frictional Unemployment
\subsection{Problem 5: Structural vs. Frictional Unemployment}
\begin{frame}{Problem 5: Structural vs. Frictional Unemployment}
In an economy, the total number of unemployed people is 8 million. Of these, 3 million are frictionally unemployed, and 2 million are structurally unemployed.
\begin{enumerate}
    \item What percentage of total unemployment is structural?  
    \item What percentage is frictional?  
    \item What portion of unemployment is due to other factors (e.g., cyclical)?  
\end{enumerate}
\end{frame}

\begin{frame}{Answer to Problem 5}
(a) Structural Unemployment Percentage:
\[
\text{Structural Unemployment} = \frac{2}{8} \times 100 = 25\%
\]

(b) Frictional Unemployment Percentage:
\[
\text{Frictional Unemployment} = \frac{3}{8} \times 100 = 37.5\%
\]

(c) Other Unemployment:
\[
\text{Other Unemployment} = 8 - (3 + 2) = 3 \, \text{million}
\]
\[
\text{Percentage} = \frac{3}{8} \times 100 = 37.5\%
\]
\end{frame}






\end{document}
